% v8: atlas definitions and branch-transition audit.
% Independent source.  No v7 file is changed.
% New bibliography keys: ShawSurfaces, JellRauShaw, LRSChern, CartwrightComplexes.
% Existing keys: EM, VuNgoc.  Existing labels: per:chart-definition,
% per:k4-realization, per:branch-duality.

\section{作用適合 atlas と分岐辺での整数アフィン延長}
\label{atl:section}

一つの cone chart と，複数の chart の tropical compatibility を区別する必要がある．
定義\ref{per:chart-definition}の同相写像は，標準 fan の整数アフィン関数の層を
引き戻すことにより，その定義域に局所的な matroidal 構造を与える．
従って，既に得た一 core の局所存在を変更する必要はない．
ただし，各 face 上の格子同型と branch の saturation だけでは，
異なる chart の間の transition が tropical isomorphism であることは従わない．

\subsection{atlas の定義と一つの欠陥整数}

本節で integer affine map とは，整数行列を線形部分に持ち，
translation は実数でよい affine map をいう．
Shaw\cite[Definition~3.1]{ShawSurfaces}，
Jell--Rau--Shaw\cite[Definitions~2.1, 5.1]{JellRauShaw}，
L\'opez de Medrano--Rinc\'on--Shaw\cite[Definitions~2.2--2.3]{LRSChern}
の chart transitions は，ambient integer affine maps の制限である．
各 face ごとに異なる affine map を選ぶだけでは，この条件を満たさない．

\begin{definition}[focus 特異点を許す作用適合 matroidal atlas]
\label{atl:definition}
$B$ を本稿の底空間，$\Delta_{\mathrm{ff}}$ を正則 faces の内部にある有限個の
focus--focus values とする．$B\setminus\Delta_{\mathrm{ff}}$ 上の
\emph{作用適合 matroidal atlas}とは，二次元 matroidal fans の開集合への
stratum-preserving homeomorphisms からなる atlas で，次を満たすものである．
各 transition とその逆は，局所的に ambient integer affine map の制限である．
各正則 face の整数余接線格子は，Lagrangian fibration の作用一形式が定める格子と一致する．
各 focus point の punctured neighborhood では，実際の focus--focus fibration の
作用アフィン構造を用いる．
\end{definition}

この定義は，focus point 自体に通常の matroidal chart があるとは要求しない．
また，punctured space に対する finite-type compactification の条件を追加していない．
その種の大域条件を要求する文献の定義を使う場合には，別途その仮定を確認する．
整数アフィン関数の層を用いれば，transition 条件は，各 chart が引き戻す層
$\operatorname{Aff}_{\ZZ}$ が overlap 上で一致することと言い換えられる
\cite[Definition~3.1]{JellRauShaw}．

三価 branch の最小 ambient lattice は rank three である．
その primitive tangent を $v$，三つの outward transverse lifts を $u_1,u_2,u_3$ とし，
\begin{equation}\label{atl:branch-relation}
 N=\ZZ v\oplus\ZZ u_1\oplus\ZZ u_2,
 \qquad u_1+u_2+u_3=\kappa v
\end{equation}
と書く．例えば product model $\RR\times\operatorname{Trop}(\text{line})$ では
$\kappa=0$ とできる．$K_4$ fan の ray $r_i$ で，残りの三本の rays を lifts に取ると
$\kappa=-1$ である．target 側にも $(v',u'_1,u'_2,u'_3,\kappa')$ を取る．

\begin{lemma}[三価 branch の延長条件]
\label{atl:branch-extension}
branch germs の同相写像 $G$ が，各 incident sheet 上で integer affine isomorphism であり，
branch tangent lattices を同型に写すとする．原点を branch 上の対応する点に移し，
sheet の対応を置換 $\sigma$ で表すと，その線形部分は
\begin{equation}\label{atl:sheet-linear}
 v\longmapsto\varepsilon v',\qquad
 u_a\longmapsto u'_{\sigma(a)}+c_a v',
 \qquad \varepsilon\in\{1,-1\},\quad c_a\in\ZZ
\end{equation}
となる．$G$ が一つの ambient affine lattice isomorphism の制限となるための
必要十分条件は
\begin{equation}\label{atl:defect-zero}
 c_1+c_2+c_3=\varepsilon\kappa-\kappa'
\end{equation}
である．その延長は一意であり，線形部分は $\operatorname{GL}(3,\ZZ)$ に属する．
\end{lemma}

\begin{proof}
branch lattice の primitive generator は $\pm v'$ に写る．
各 sheet の格子を branch lattice で割ると rank-one lattice となる．
outward half-plane を outward half-plane へ写す unimodular map は，
その正の primitive generator を正の primitive generator に写すので，
\eqref{atl:sheet-linear}を得る．三 sheet の translations は branch 上で一致し，
基点を移せば一つの translation を除去できる．

三本の $u_a$ の像の和は $(\kappa'+\sum_a c_a)v'$ であり，
source の relation の像は $\varepsilon\kappa v'$ である．
従って\eqref{atl:defect-zero}が必要である．逆にこの式が成り立てば，
$(v,u_1,u_2)$ の像を指定する線形写像は $u_3$ 上でも所定の写像となる．
$(v',u'_{\sigma(1)},u'_{\sigma(2)})$ も整数基底であるから，
この写像の determinant は符号を除いて一である．
star は ambient space を張るので延長は一意である．
\end{proof}

\begin{example}[facewise unimodularity だけでは十分でない]
\label{atl:counterexample}
$v=e_0$，$u_1=e_1$，$u_2=e_2$，$u_3=-e_1-e_2$ とする．
第1 sheet のみで
\[
 xv+tu_1\longmapsto(x+t)v+tu_1\qquad(t\geq0)
\]
とし，第2・第3 sheet では恒等写像とする．
これは branch を点ごとに固定する同相写像で，全 sheet で unimodular であり，
branch lattices も飽和する．しかし $\sum_a c_a=1$ なので
ambient integer affine extension は存在しない．
これは二つの局所 chart の非適合性の例であり，一つの chart が
matroidal 構造を引き戻せないという意味ではない．
\end{example}

\subsection{共通の Hamiltonian による十分条件}

\begin{proposition}[共通 circle action による延長]
\label{atl:common-hamiltonian}
補題\ref{atl:branch-extension}の状況で，branch の近傍に連続関数 $h$ があり，
source と target の両 chart において，$h$ が一つの ambient affine function の
pullback として表されるとする．$h$ の branch tangent に沿う勾配が非零ならば，
$G$ は ambient affine lattice isomorphism へ一意に延長する．
特に，$h$ を両側で同じ共通 Hamiltonian circle の作用積分として取れればよい．
二つの primitive の違いによる定数の差は許す．
\end{proposition}

\begin{proof}
基点と定数項を除き，$h$ の source と target における線形部分を
$\ell$，$\ell'$ と書く．各 sheet で $\ell=\ell'\circ G$ が成り立つ．
三つの transverse lifts についてこの式を足すと
\[
 \kappa\ell(v)
 =\sum_a\ell(u_a)
 =\ell'\left(\sum_a u'_{\sigma(a)}+\sum_a c_a v'\right)
 =(\kappa'+\sum_a c_a)\ell'(v').
\]
一方 $\ell(v)=\varepsilon\ell'(v')$ であり，$\ell'(v')\ne0$ である．
従って\eqref{atl:defect-zero}を得る．補題\ref{atl:branch-extension}を適用する．
\end{proof}

ここで必要なのは，共通 circle の一つの Hamiltonian の ambient-affine 性である．
block 全体の $H_1$ のすべての元に ambient affine function を割り当てる条件ではない．
$K_4$ chart では\eqref{per:k4-lattice-map}の $\Phi(\mu_i)$ がこの関数を与える．
Petersen chart でも，補題\ref{per:branch-duality}の
$\alpha\mapsto(-1,0,0)$ が各 branch でこの条件を満たす．

\begin{lemma}[denominator one の branch chart]
\label{atl:qone-chart}
三価 branch block の共通 circle action を $h$ とする．
各 sheet に作用座標 $(h,J_a)$ があり，branch へ連続に延長すると仮定する．
critical action curves が
\[
 J_a=m_a h+b_a,\qquad m_a\in\ZZ,\quad b_a\in\RR
\]
という直線であり，その両側の regular sheets で
$\bigl(h,\delta_a(J_a-m_a h-b_a)\bigr)$ が branch までの half-plane の
局所同相座標となるとする．ここで $\delta_a=\pm1$ は outward direction を選ぶ符号である．
このとき，共通 Hamiltonian $h$ が ambient-affine となる作用適合 branch chart が存在する．
\end{lemma}

\begin{proof}
$\kappa=0$ の標準 star を取り，各 sheet 上で
\[
 F_a=h v+\delta_a(J_a-m_a h-b_a)u_a
\]
と定める．各 face map の作用座標に関する determinant は $\delta_a$ であり，
branch ではすべて $F_a=hv$ に一致する．仮定した half-plane 座標により
貼り合わせた $F$ は star の開近傍への同相写像である．
$h$ は $v$ に一，$u_1,u_2$ に零を与える一つの整数線形形式の pullback である．
\end{proof}

simple saddle からの通常の action coordinates では，最後の half-plane 条件は
transverse action の連続性と strict monotonicity によって確認できる．
これに対し，critical action curve の直線性は，Lagrangian 葉層の補間からは従わない．
その条件は実際の parametric area calculation で証明する必要がある．

\subsection{大域貼り合わせの条件と scope}

\begin{proposition}[局所作用座標から atlas への貼り合わせ]
\label{atl:gluing-criterion}
$B\setminus\Delta_{\mathrm{ff}}$ を，既に構成された一 core の matroidal charts，
三価 branch charts，および正則 face の action charts で覆う．
core charts の core points を含む部分は互いに交わらないように縮めてよい．
次を仮定する．
各 chart は同じ actual fibration の正則 face の作用格子に適合する．
branch に交わる任意の overlap で，その共通 circle action の作用積分は，
両 chart で一つの ambient affine function として表され，branch 上で非定数である．
このとき，これらの charts は定義\ref{atl:definition}の作用適合 matroidal atlas を与える．
\end{proposition}

\begin{proof}
正則 face 内の overlap では，二組の作用座標は整数基底の変換と実数 translation で異なる．
face lattice は ambient lattice の飽和部分格子なので，必要なら整数補空間を選ぶことにより，
この rank-two affine map は ambient integer affine map として表せる．
branch を含む overlap では命題\ref{atl:common-hamiltonian}を用いる．
branch の rank-three span がより大きな core ambient lattice に飽和して入っている場合も，
整数補空間を選べば同じ結論を得る．
最後に transition は実際の chart maps の合成 $F_b\circ F_a^{-1}$ なので，
triple overlap 上の cocycle identity は写像として自動的に成立する．
\end{proof}

cocycle identity が要求されるのは実際の overlap 上である．
overlap が一つの二次元 face にのみ含まれる場合，ambient rank-three space 全体で
選んだ延長行列の積まで恒等になる必要はない．
そのような延長は face の外では一意ではないからである．
これに対し，branch star を含む overlap では ambient span が rank three なので，
補題\ref{atl:branch-extension}の一意性も使える．

Evans--Mauri\cite[Remark~2.7, Lemma~3.1 and \S3.1]{EM}は，正則部の
作用アフィン構造と，四直線補集合から tetrahedral-link cone への
Lagrangian fibration を与える．同論文の Proposition~6.13 と Theorem~6.15 の
大域 monodromy 比較は，Assumption~6.6 の log Calabi--Yau 条件を持つ
one-dimensional strata の toric neighborhoods に対する結果である．
三 punctures を持つ line と branch sheets を含む本節の貼り合わせを，
その定理から直接結論することはできない．

Cartwright の tropical complex は，dual complex に intersection-number data
$\alpha(v,e)$ を加え，codimension-one star 上の線形関数を別途指定する構造である
\cite[Definitions~2.5, 2.8; Construction~3.2]{CartwrightComplexes}．
その sheaf の指定は，本節で facewise affine maps と ambient affine maps を
区別する必要性とよく対応する．ただし，tropical complex の存在だけでは，
matroidal smoothness や本稿の actual action periods との一致は従わない．
特異 affine 構造の置かれる場所が異なり得る点についても，
同論文 Remark~3.4 が Gross--Siebert 型構造との違いを明示している．

本稿の ordered small resolution では，この違いを intersection matrix から
直接確認できる．次の命題は作用 atlas の不存在を述べるものではない．

\begin{proposition}[ordered degeneration の weak tropical complex]
\label{atl:cartwright-distinction}
次数分割 $d=e_1+\cdots+e_k$ の ordered small resolution から得る
Cartwright の weak tropical complex は，最初の component $D_1$ において
tropical complex の local intersection condition を満たさない．
具体的に，$\mathbf e=(e_2,\ldots,e_k)^{\mathsf T}$ と書くと，その local matrix は
\begin{equation}\label{atl:negative-cartwright}
 Q_1=e_1\left(\mathbf e\mathbf e^{\mathsf T}
                  -d\,\operatorname{diag}(e_2,\ldots,e_k)\right)
\end{equation}
であり，負定値である．全平面分割では $Q_1=J_{d-1}-dI_{d-1}$ となる．
\end{proposition}

\begin{proof}
$D_1$ 上の boundary curves は
$C_{1j}=e_jH_1-\sum_{p=1}^{d e_1e_j}E_{1j,p}$ である．
従って
\[
 C_{1j}^2=e_1e_j(e_j-d),\qquad
 C_{1j}\cdot C_{1l}=e_1e_je_l\quad(j\ne l),
\]
すなわち\eqref{atl:negative-cartwright}を得る．
Cartwright の convention では $\alpha(v,e)$ は，$e$ の反対側の
surface 内の double curve の自己交点の負である
\cite[Lemma~2.2]{CartwrightComplexes}．
従って，Definition~2.8 の $D_1$ における local intersection matrix は正確に $Q_1$ である．
任意の非零 $z=(z_2,\ldots,z_k)$ に対し，weighted Cauchy--Schwarz inequality より
\[
 \frac1{e_1}z^{\mathsf T}Q_1z
 =\left(\sum_{j=2}^ke_jz_j\right)^2-d\sum_{j=2}^ke_jz_j^2
 \leq-e_1\sum_{j=2}^ke_jz_j^2<0.
\]
local matrix に正の固有値がちょうど一つあるという条件に反する．
全平面分割では固有値は $-1$ と，重複度 $d-2$ の $-d$ である．
\end{proof}

したがって，本稿で考える matroidal action atlas を，元の semistable dual complex に
Cartwright の intersection data を付けただけの構造と同一視してはならない．
node neighborhoods を別の strata として残すことや，作用周期から chart を構成することは，
この比較においても実質的な追加データである．

最後に，simple focus--focus point の周りの作用格子は
$\left(\begin{smallmatrix}1&1\\0&1\end{smallmatrix}\right)$ と共役な非自明 monodromy を持つ
\cite[Remark~3.1]{VuNgoc}．その点が位相的に face の内部にある以上，
そこに通常の smooth matroidal chart を延長すれば局所モデルは $\RR^2$ となり，
punctured disk の monodromy は自明になるはずである．従って，
focus points の除外または明示的な affine singularities の許容は不可欠である．
